%% =====================================================================
%%  metadata.tex  --  thông tin bài báo in trên hai trang bìa
%%  Sửa tiêu đề, tác giả, đơn vị, tóm tắt, từ khóa ở đây.
%% =====================================================================

%% ---------------------------------------------------------------------
%% 1. TIÊU ĐỀ
%% ---------------------------------------------------------------------
\jsttitlevi{%
  Lập đường trên bề mặt dẫn đường thành phố ba chiều bằng Funnel Tree
  và kiểm chứng mô phỏng UAV%
}
\jsttitleen{%
  Surface-Constrained Route Planning on a 3D City Navigation Surface
  Using Funnel Trees and UAV Simulation Validation%
}

%% ---------------------------------------------------------------------
%% 2. LOẠI BÀI (in ở góc phải header)
%% ---------------------------------------------------------------------
\jstarticletypevi{Bài nghiên cứu}
\jstarticletypeen{Research Article}

%% ---------------------------------------------------------------------
%% 3. TÁC GIẢ
%%    Cú pháp: \jstauthorvi[chỉ số đơn vị]{Họ và tên}
%%    - Chỉ số đơn vị in thành superscript, khớp với \jstaffiliationvi.
%%    - Dấu * đánh dấu tác giả liên hệ (khai báo thêm ở mục 5).
%% ---------------------------------------------------------------------
\jstauthorvi[1,*]{Phan Thành An và các tác giả khác}
\jstauthoren[1,*]{Phan Thanh An and other authors}

%% ---------------------------------------------------------------------
%% 4. ĐƠN VỊ
%% ---------------------------------------------------------------------
\jstaffiliationvi{1}{%
  Viện Toán học và Các khoa học tính toán (IMACS), Trường Đại học Bách khoa
  TP. Hồ Chí Minh (HCMUT), 268 Lý Thường Kiệt, Phường Diên Hồng,
  TP. Hồ Chí Minh, Việt Nam.%
}
\jstaffiliationen{1}{%
  Institute of Mathematical and Computational Sciences (IMACS),
  Ho Chi Minh City University of Technology (HCMUT), 268 Ly Thuong Kiet,
  Dien Hong Ward, Ho Chi Minh City, Vietnam.%
}

%% ---------------------------------------------------------------------
%% 5. TÁC GIẢ LIÊN HỆ  {Họ tên}{Địa chỉ}{Email}
%% ---------------------------------------------------------------------
\jstcorrespondingvi{Phan Thành An}{%
  Viện Toán học và Các khoa học tính toán (IMACS), Trường Đại học Bách khoa
  TP. Hồ Chí Minh (HCMUT), 268 Lý Thường Kiệt, Phường Diên Hồng,
  TP. Hồ Chí Minh, Việt Nam.}{thanhan@hcmut.edu.vn}

\jstcorrespondingen{Phan Thanh An}{%
  Institute of Mathematical and Computational Sciences (IMACS),
  Ho Chi Minh City University of Technology (HCMUT), 268 Ly Thuong Kiet,
  Dien Hong Ward, Ho Chi Minh City, Vietnam.}{thanhan@hcmut.edu.vn}

%% ---------------------------------------------------------------------
%% 6. TÓM TẮT VÀ TỪ KHÓA
%% ---------------------------------------------------------------------
\jstabstractvi{%
Nghiên cứu trình bày một quy trình lập đường và kiểm chứng mô phỏng cho thiết
bị bay không người lái (Unmanned Aerial Vehicle, UAV) trong môi trường thành
phố ba chiều. Bản triển khai Python của thuật toán Funnel Tree được kiểm chứng
trước hết trên chín lưới đa diện chuẩn, gồm bề mặt kín, bề mặt không lồi, mặt
xuyến và hai bề mặt mở; năm lưới có kết quả tham chiếu với sai số tương đối không
vượt $10^{-4}$; bốn lưới còn lại được dùng để kiểm tra khả năng chạy trên bề mặt
mở, bề mặt lồi và bề mặt có lỗ, không có reference độc lập. Quy trình sau đó được áp dụng cho mô hình thành phố
\textit{Ho Chi Minh City 2025} (\num{301490} tam giác). Ba bề mặt dẫn đường được
dựng với ô danh nghĩa $5$, $4$ và $3$ m, lần lượt có \num{2295}, \num{3355} và
\num{5913} đỉnh; mỗi bề mặt là envelope an toàn phía trên công trình, có khoảng
hở đứng, dự phòng ngang và vùng chuyển tiếp cao độ. Trên lưới $4$ m, ba tuyến
từ một nguồn đến ba đích dài \num{142.909050}, \num{127.295935} và
\num{86.194146} mét của mô hình thu nhỏ. Các tuyến được tham số hóa bằng spline
Catmull--Rom và một profile vận tốc thang, rồi bay trong Gazebo Jetty 10.5 bằng
mô hình khối điểm $1$ kg với điều khiển bù trước và phản hồi PD. Chín chuyến bay
(ba lưới $\times$ ba đích) đều hợp lệ; bảy chuyến đạt mốc sai số bám $1$ m, trong
đó lưới $4$ m là cấu hình duy nhất đạt cả ba với sai số lớn nhất
$0{,}399$--$0{,}601$ m. Lưới $5$ m và $3$ m mỗi lưới có một ca chưa đạt ngưỡng sai số
bám $1$ m ở đích thứ hai (lần lượt $1{,}022$ m và $1{,}381$ m).
Các ca này xuất hiện tại góc cua có waypoint phân bố không đều
và vùng chuyển tiếp cao độ dốc $45^\circ$, với cấu hình vận tốc
tham chiếu tối đa $5$ m/s và gia tốc tiếp tuyến tối đa $3$ m/s$^2$.
Kết quả xác nhận khả năng hoạt động của quy trình trong phạm vi ba
cấu hình và ba đích; nghiên cứu chưa giải quyết tối ưu đường bay
trong toàn bộ không gian tự do ba chiều hoặc điều phối nhiều UAV.%
}

\jstkeywordsvi{%
Funnel Tree; đường đi ngắn nhất; bề mặt đa diện; thành phố 3D; UAV; Gazebo;
trắc địa bề mặt%
}

\jstabstracten{%
This study presents a route-planning and simulation-validation workflow for a
single unmanned aerial vehicle in a three-dimensional city environment. The
Python implementation of the Funnel Tree algorithm is first verified on nine
benchmark polyhedral meshes, covering closed surfaces, a non-convex surface, a
toroidal surface and two open surfaces; five meshes match the reference results
within a relative error of $10^{-4}$; the remaining meshes exercise open, convex and toroidal
surfaces and have no independent reference. The workflow is then applied to
the \textit{Ho Chi Minh City 2025} city model (\num{301490} triangles). Three
navigation surfaces are built with nominal cell sizes of $5$, $4$ and $3$ m,
with \num{2295}, \num{3355} and \num{5913} vertices, respectively; each is a
safety envelope above the buildings with a vertical clearance, a horizontal
margin and elevation transitions. On the $4$ m grid, three routes from one
source to three destinations are 142.909050, 127.295935 and 86.194146 metres of
the scaled model long. The routes are parameterized by a
Catmull--Rom spline and a single trapezoidal speed profile, and flown in Gazebo
Jetty 10.5 with a $1$ kg point-mass model under feedforward and PD control. All
nine flights (three grids $\times$ three goals) are valid; seven meet the $1$ m
tracking target, and the $4$ m grid is the only configuration that meets it on
all three goals, with maximum errors of 0.399--0.601 m. The $5$ m and $3$ m grids each have one case that does not meet
the $1$ m tracking-error target at the second goal
(1.022 m and 1.381 m, respectively). These cases occur at a corner
with uneven waypoint spacing and a $45^\circ$ elevation transition,
under a reference speed limit of $5$ m/s and a tangential
acceleration limit of $3$ m/s$^2$. The results
support the feasibility of the integrated workflow for the three tested
configurations and three goals; the study does not establish optimality in
unrestricted three-dimensional free space or demonstrate real-UAV operation or
multi-UAV coordination.%
}

\jstkeywordsen{%
Funnel Tree; shortest path; polyhedral surface; 3D city; UAV; Gazebo; surface
geodesic%
}

%% ---------------------------------------------------------------------
%% 7. KHUNG "TRÍCH DẪN BÀI BÁO NÀY" Ở ĐÁY BÌA
%% ---------------------------------------------------------------------
\jstcitationvi{%
An PT và cs. Lập đường trên bề mặt dẫn đường thành phố ba chiều bằng Funnel
Tree và kiểm chứng mô phỏng UAV. Tạp chí Khoa học và Công nghệ ĐHQG-HCM,
chuyên san Kỹ thuật và Công nghệ. 2026.%
}

\jstcitationen{%
An PT, et al. Surface-Constrained Route Planning on a 3D City Navigation
Surface Using Funnel Trees and UAV Simulation Validation. VNUHCM Journal of
Engineering and Technology. 2026.%
}
